\chapter{Modelos de Referências}
\label{chap:modelos-referencias}

Este capítulo reúne estruturas de conferência da ABNT NBR 6023:2025. Campos entre colchetes angulares devem ser substituídos pelos dados reais da fonte. Não mantenha marcadores no trabalho final e não complete lacunas por suposição.

\caixaatencao{Consulte a própria fonte, a página editorial, o catálogo ou a base responsável. Geradores automáticos podem ajudar, mas cada referência precisa ser conferida. Os endereços terminados em ``.exemplo'' e o DOI com ``10.xxxx'' são marcadores didáticos visíveis e não constituem hiperlinks.}

\section{Regras comuns}

As referências são alinhadas à margem esquerda, em espaço simples e separadas entre si por uma linha em branco. O recurso tipográfico usado para destacar o título deve ser uniforme. Autoria, título, edição, local, editora e data são registrados conforme o tipo documental e os dados disponíveis.

\section{Livro impresso}

\refmodelo{SOBRENOME, Prenome. \textbf{Título}: subtítulo. Número da edição, a partir da segunda. Local: Editora, ano.}

Registre responsabilidade secundária, volume, coleção ou tradução quando forem relevantes para identificar a edição consultada.

\section{Livro eletrônico}

\refmodelo{SOBRENOME, Prenome. \textbf{Título}: subtítulo. Edição. Local: Editora, ano. Formato ou suporte, quando necessário. Disponível em: \urlexemplo{https://endereco.exemplo/livro}. Acesso em: dia mês abreviado. ano.}

Não use o endereço de uma página de resultados quando houver página estável do livro. Se houver DOI, registre-o conforme aparece na fonte e mantenha os dados exigidos para o documento consultado.

\section{Capítulo de livro}

\refmodelo{SOBRENOME, Prenome do autor da parte. Título do capítulo. In: SOBRENOME, Prenome do responsável pela obra (org.). \textbf{Título do livro}: subtítulo. Edição. Local: Editora, ano. páginas inicial--final.}

A autoria do capítulo e a responsabilidade pelo livro são informações diferentes. Não atribua o capítulo ao organizador quando houver autoria própria.

\section{Artigo de periódico sem DOI}

\refmodelo{SOBRENOME, Prenome. Título do artigo. \textbf{Título do periódico}, local, volume, número, páginas ou identificador eletrônico, período e ano. Disponível em: \urlexemplo{https://endereco.exemplo/artigo}. Acesso em: dia mês abreviado. ano.}

Em artigo impresso, omitem-se endereço e acesso. Use o identificador eletrônico no lugar das páginas quando essa for a forma adotada pelo periódico.

\section{Artigo de periódico com DOI}

\refmodelo{SOBRENOME, Prenome. Título do artigo. \textbf{Título do periódico}, local, volume, número, páginas ou identificador, ano. DOI: \urlexemplo{https://doi.org/10.xxxx/xxxxx}. Disponível em: \urlexemplo{https://endereco-consultado.exemplo}. Acesso em: dia mês abreviado. ano.}

O DOI é um identificador persistente; a URL informa onde o conteúdo foi consultado. Não invente DOI e não use a expressão ``DOI ou Disponível em'' como se os elementos fossem sempre intercambiáveis. Quando a consulta for eletrônica, registre os elementos exigidos para acesso conforme a norma e o tipo documental.

\section{TCC, dissertação e tese}

\refmodelo{SOBRENOME, Prenome. \textbf{Título}: subtítulo. Ano de depósito. Tipo do trabalho (Grau e área) -- Unidade acadêmica, Instituição, Local, ano de defesa.}

Em versão eletrônica, acrescente o endereço permanente e a data de acesso. Diferencie o ano de depósito do ano de defesa quando a fonte apresentar ambos.

\section{Trabalho em evento}

\refmodelo{SOBRENOME, Prenome. Título do trabalho. In: NOME DO EVENTO, número., ano, local de realização. \textbf{Título dos anais [...]} Local de publicação: Responsável pela publicação, ano. páginas ou identificador. Disponível em: \urlexemplo{https://endereco.exemplo}. Acesso em: dia mês abreviado. ano.}

Não confunda local do evento com local de publicação. Para evento on-line, registre os dados efetivamente apresentados pela fonte.

\section{Legislação}

\refmodelo{BRASIL. Lei nº 13.709, de 14 de agosto de 2018. Lei Geral de Proteção de Dados Pessoais. \textbf{Diário Oficial da União}: Brasília, DF, 15 ago. 2018. Disponível em: \url{https://www.planalto.gov.br/ccivil_03/_ato2015-2018/2018/lei/l13709.htm}. Acesso em: 6 ago. 2026.}

Identifique a jurisdição, o tipo e número do ato, a data, a ementa ou título, a publicação oficial e os dados de acesso. Consulte o texto consolidado quando houver alterações.

\section{Jurisprudência}

\refmodelo{JURISDIÇÃO. Órgão julgador. Tipo e número do processo. Relator: [nome]. Data do julgamento. \textbf{Fonte de publicação}, local, data. Disponível em: \urlexemplo{https://endereco-oficial.exemplo}. Acesso em: dia mês abreviado. ano.}

Registre os elementos capazes de identificar a decisão específica. Não referencie apenas a página inicial do tribunal.

\section{Ato administrativo}

\refmodelo{DISTRITO FEDERAL. Universidade do Distrito Federal Professor Jorge Amaury Maia Nunes. Instrução nº 24, de 15 de abril de 2025. Aprova o Regimento da Graduação da UnDF. \textbf{Diário Oficial do Distrito Federal}: Brasília, DF, nº 74, p. 10, 22 abr. 2025.}

Portarias, instruções, resoluções e editais devem identificar a entidade, o ato, a data, a ementa e a fonte oficial. Acrescente endereço e acesso quando consultados on-line.

\section{Documento institucional}

\refmodelo{ENTIDADE RESPONSÁVEL. \textbf{Título do documento}: subtítulo. Versão ou edição. Local: Entidade, ano. Disponível em: \urlexemplo{https://endereco-institucional.exemplo}. Acesso em: dia mês abreviado. ano.}

Use a instituição como autora somente quando ela assumir responsabilidade intelectual pelo documento.

\section{Relatório técnico ou científico}

\refmodelo{AUTOR OU ENTIDADE. \textbf{Título do relatório}: subtítulo. Local: Instituição responsável, ano. Número ou série do relatório. Disponível em: \urlexemplo{https://endereco.exemplo}. Acesso em: dia mês abreviado. ano.}

Inclua número, contrato, projeto ou série quando forem necessários para distinguir o relatório.

\section{Norma técnica}

\refmodelo{ASSOCIAÇÃO BRASILEIRA DE NORMAS TÉCNICAS. \textbf{ABNT NBR 10520: informação e documentação -- citações em documentos -- apresentação}. Rio de Janeiro: ABNT, 2023.}

Use o código, o título e o ano da edição efetivamente consultada. Não cite como vigente uma edição substituída.

\section{Matéria jornalística}

\refmodelo{SOBRENOME, Prenome. Título da matéria. \textbf{Nome do jornal}, local, dia mês abreviado. ano. Seção, página.}

Na versão eletrônica, acrescente endereço e acesso. Quando não houver autoria pessoal, a entrada pode começar pelo título, conforme a regra aplicável.

\section{Vídeo}

\refmodelo{TÍTULO do vídeo. Responsável principal: [nome]. [Local]: Produtora ou canal, ano. 1 vídeo (duração). Disponível em: \urlexemplo{https://endereco.exemplo}. Acesso em: dia mês abreviado. ano.}

Registre a responsabilidade relevante, a plataforma ou produtora, a data e a duração. Use timestamp na citação direta de trecho audiovisual.

\section{Podcast}

\refmodelo{TÍTULO do episódio. In: \textbf{Título do podcast}. Apresentação: [nome]. [Local]: Produtora, dia mês abreviado. ano. Podcast (duração). Disponível em: \urlexemplo{https://endereco.exemplo}. Acesso em: dia mês abreviado. ano.}

Distinga o título do episódio do título da série e identifique apresentador, entrevistado ou produtor conforme a responsabilidade efetiva.

\section{Entrevista}

\refmodelo{SOBRENOME, Prenome do entrevistado. \textbf{Título ou assunto}. Entrevista concedida a [nome do entrevistador]. Local, data. Suporte, duração ou localização, quando necessário.}

Entrevistas produzidas pela pesquisa também precisam ser descritas metodologicamente. Use código do participante quando a identificação for vedada e preserve os registros segundo o protocolo ético.

\section{Imagem, fotografia ou ilustração}

\refmodelo{AUTOR OU ENTIDADE. \textbf{Título}. Ano. 1 imagem, técnica ou tipo, dimensões quando relevantes. Acervo ou plataforma. Disponível em: \urlexemplo{https://endereco.exemplo}. Acesso em: dia mês abreviado. ano.}

Além da referência, verifique autoria, licença e autorização de uso. ``Fonte: Google'' não identifica a obra nem concede licença.

\section{Mapa ou documento cartográfico}

\refmodelo{AUTOR OU ENTIDADE. \textbf{Título}: subtítulo. Local: Responsável, ano. Escala. Tipo e suporte. Disponível em: \urlexemplo{https://endereco.exemplo}. Acesso em: dia mês abreviado. ano.}

Registre escala, projeção, folha, série e coordenadas quando forem relevantes para identificar ou reproduzir a análise.

\section{Página da internet}

\refmodelo{AUTOR OU ENTIDADE. \textbf{Título da página}. Local ou portal, data de publicação ou atualização. Disponível em: \urlexemplo{https://endereco-especifico.exemplo}. Acesso em: dia mês abreviado. ano.}

Prefira a página específica. Não use apenas o nome do buscador, a página inicial do portal ou um endereço encurtado se houver URL estável.

\section{Base de dados}

\refmodelo{ENTIDADE RESPONSÁVEL. \textbf{Nome da base de dados}. Versão ou período. Local: Entidade, ano. Disponível em: \urlexemplo{https://endereco.exemplo}. Acesso em: dia mês abreviado. ano.}

Informe versão, recorte, tabela, painel, filtros e data de extração no texto ou na metodologia quando os resultados puderem mudar.

\section{Conjunto de dados}

\refmodelo{AUTOR OU ENTIDADE. \textbf{Título do conjunto de dados}. Versão. Repositório, ano. Conjunto de dados. DOI: \urlexemplo{https://doi.org/10.xxxx/xxxxx}. Disponível em: \urlexemplo{https://endereco.exemplo}. Acesso em: dia mês abreviado. ano.}

Indique versão, responsável, repositório, identificador persistente e data de acesso. Não confunda a referência do conjunto com a do artigo que o analisou.

\section{Publicação em rede social}

\refmodelo{AUTOR OU PERFIL RESPONSÁVEL. Início ou título da publicação. Local: Plataforma, dia mês abreviado. ano. Tipo de publicação. Disponível em: \urlexemplo{https://endereco-direto.exemplo}. Acesso em: dia mês abreviado. ano.}

Use redes sociais somente quando forem fonte pertinente ao objeto. Preserve evidência da publicação quando a mutabilidade for relevante e observe ética, privacidade e termos de uso.

\section{Documento sem autoria pessoal}

Procure primeiro autoria institucional. Se nenhuma responsabilidade puder ser identificada, a entrada começa pelo título conforme a NBR 6023:2025. A ausência de autoria exige avaliação crítica reforçada e não autoriza atribuir o conteúdo a uma instituição apenas porque ela hospeda o arquivo.

\section{Local, editora ou data não identificados}

Não omita silenciosamente elementos nem invente dados. Use informação inferida entre colchetes quando houver base segura. Se o local não puder ser identificado, empregue [S. l.]; se a editora não puder ser identificada, [s. n.]; se ambos faltarem, [S. l.: s. n.]. Para data ausente, indique uma data provável ou intervalo entre colchetes, como [2024?] ou [entre 2020 e 2024], somente quando a inferência puder ser justificada.

\section{Conferência final}

Antes da entrega, confira:

\begin{itemize}
  \item autoria e responsabilidade conforme a fonte;
  \item título, edição, local, editora e data;
  \item DOI, URL estável e data de acesso quando aplicáveis;
  \item funcionamento dos links e correspondência com o documento citado;
  \item uniformidade de destaque, pontuação e abreviaturas;
  \item presença de toda fonte citada na lista;
  \item uso efetivo de toda fonte incluída na lista.
\end{itemize}
